In inertial pavement screening, gains in the selected classifier's outer score can accompany less accurate validation estimates or unchanged coverage of annotated road events. We examine these outcomes using fixed-candidate nested selection on nine PVS acquisitions and an independent Karlsruhe date test with spatial purging and matched inner training budgets. RoadSens source and temporal controls establish the evaluation context. In PVS disjoint windows, matched selection raises scenario outer macro-F1 from 0.7904 to 0.8168 while estimation MAE rises from 0.1069 to 0.1290; the primary two-axis criterion is unmet. On the Karlsruhe final date, buffered-date selection raises macro-F1 from 0.6782 to 0.7855 and reduces absolute estimation error from 0.2817 to 0.2072, while both systems detect 7/12 pothole intervals. PVS scenario interval coverage rises from 34/60 to 52/60 alongside more non-bump alarms. One-second inertial statistics improve all eight matched RoadSens comparisons. The results distinguish classifier selection, score estimation and event coverage, with acquisition groups, refit budgets and alarm exposure attached to each screening claim.
